\subsection{Divisibility on a total-product divisor}
\label{subsec:degree-three-saturation}

We strengthen the coefficient injection by computing
its quotient along a total-product divisor. This controls
the residues of the odd abelian components that first
reappear in the class-three extension.
Take a finite extension $L/\mathbb Q_p$ splitting the
characteristic polynomial on $(T/\pi Z)[1/p]$ and write
$\mathcal S_L=\mathcal S\otimes_{\mathbb Z_p}L$.
Tensor coefficients below have a common bounded
denominator in $L$. Put $z_i=1+X_i$ and
\[
 \mathcal A_3=O_L[[X_1,X_2,X_3]][1/p],\qquad
 \mathcal B_3=\mathcal A_3
       [q(z_1z_2)^{-1},q(z_1z_3)^{-1},q(z_2z_3)^{-1}],
\]
where $q$ is a characteristic polynomial for the finite
rational module $(T/\pi Z)[1/p]$.
The preceding proof constructs an equivariant injection
$\iota_3:\mathcal S_L\to\mathcal B_3^{\oplus d}$,
with $d$ the finite number of ordered color coordinates.

\begin{proposition}\label{prop:degree-three-saturation}
For $\beta\in1+\mathfrak m_L$ and every $k\geq1$,
\[
 \iota_3(\mathcal S_L)\cap
 (z_1z_2z_3-\beta)^k\mathcal B_3^{\oplus d}
       =\iota_3((\gamma-\beta)^k\mathcal S_L).
\]
Thus the induced map on $(\gamma-\beta)$-coinvariants
is injective.
\end{proposition}
\begin{proof}
We first give the exact evaluation map that keeps
divisibility inside the source. For a compact
$O_L$-module $E$ with continuous pro-$p$ action of
$\Gamma$, set $A=\beta\gamma_E^{-1}-1$.
This operator tends to zero in powers on every finite
invariant quotient: modulo the maximal ideal,
$\gamma_E-1$ is nilpotent for the finite $p$-group action.
Consequently the evaluation map is convergent and gives
\begin{equation}\label{eq:total-evaluation}
 \frac{O_L[[X_1]]\widehat\otimes E}
 {(z_1\gamma_E-\beta)(O_L[[X_1]]\widehat\otimes E)}
 \xrightarrow{\ \sim\ }E,\qquad
        \sum_{n\geq0}e_nX_1^n\longmapsto\sum_{n\geq0}A^ne_n.
\end{equation}
Here is the exact kernel calculation. Subtract the value
of the series and divide each $X_1^n-A^n$ by $X_1-A$.
The coefficient of $X_1^j$ in the quotient is the
convergent sum $\sum_{n\geq j+1}A^{n-1-j}e_n$.
This constructs a preimage in the same compact module.
Multiplication by the automorphism $\gamma_E$ identifies
the ideal generated by $X_1-A$ with that in
\eqref{eq:total-evaluation}.
The construction remains valid after inverting $p$.

Over $O_L[[X]][1/p]$, the module $W_L$ is finite free.
The bounded lattice comparison in Section~\ref{sec:abelian}
therefore permits a finite direct sum of
\eqref{eq:total-evaluation} for the first-slot colors.
Apply it both to $N$ and to $\bar E$.
The inclusion $\bar E_L\hookrightarrow N_L[1/q(\gamma)]$
is $e\mapsto q(\gamma)^{-1}q(\gamma)e$ and is compatible
with the evaluation maps. On coefficients the evaluation is
\[
                         z_1=\frac{\beta}{z_2z_3}.
\]
It follows in particular that a class from
$(W\widehat\otimes\bar E)_L$ whose evaluated tensor
lies in the regular $N_L$ coefficient module has a
representative in $(W\widehat\otimes N)_L$.
This is exactly the intersection $\bar E_L\cap N_L=N_L$
inside $N_L[1/q(\gamma)]$, transported by
\eqref{eq:total-evaluation}.

Put $D_\beta=z_1z_2z_3-\beta$.
The quotient $\mathcal A_3/(D_\beta)$ is
$O_L[[X_1,X_2]][1/p]$, by the integral change of
variable $z_3=\beta/(z_1z_2)$. For a root $\alpha$ of $q$,
the three pair divisors reduce, up to units, to
\[
 z_1-\beta/\alpha,\qquad z_2-\beta/\alpha,\qquad
                                  z_1z_2-\alpha.
\]
All roots lie in $1+\mathfrak m_L$ by continuity of the
pro-$p$ action on the finite free lattice.
These are prime factors in the bounded power-series
ring, and the three displayed families have no common
prime factors with one another. For example substituting
$z_1=\beta/\alpha'$ into the third gives a nonzero linear
polynomial in the independent variable $z_2$.
The ring is factorial, being a localization of the
regular local ring $O_L[[X_1,X_2]]$.
The intersection of its localizations in the separate
three factor families is therefore the unlocalized ring.

Represent $s\in\mathcal S_L$ by
$F\in(W\widehat\otimes\bar E)_L$.
Its tensor is alternating in the last two slots and
has poles only along $q(z_2z_3)$.
The bracket map is $F\mapsto F-\tau F$, where
$\tau$ is the cyclic permutation of three slots.
Suppose that $\iota_3(s)$ is divisible by $D_\beta$.
Modulo $D_\beta$, this says $\bar F=\tau\bar F$.
Thus $\bar F$ belongs to all three separate pair-divisor
localizations and is regular by the intersection just proved.
Equation~\eqref{eq:total-evaluation} supplies
$G\in(W\widehat\otimes N)_L$ such that
\[
 F-G\in(\gamma-\beta)(W\widehat\otimes\bar E)_L.
\]
The reduced tensor $\bar G$ is cyclically invariant and
alternating in its last two slots, hence fully alternating.
Apply the alternating idempotent in $L[S_3]$ to $G$.
The total-product divisor is symmetric in all three
positions, so the resulting alternating tensor has the
same reduction. Its denominator $6$ is one fixed rational
scalar and preserves boundedness. Fully alternating
tensors are exactly the images of
$J(\widehat{\bigwedge}^3W)_L$ in
\eqref{eq:class3-S}. Consequently
\[
 F\in J(\widehat{\bigwedge}^3W)_L
             +(\gamma-\beta)(W\widehat\otimes\bar E)_L,
\]
which proves $s\in(\gamma-\beta)\mathcal S_L$.
The converse is equivariance. Cancellation of one
$D_\beta$ in the domain $\mathcal B_3$, followed by
iteration, proves the equality for every $k$.
\end{proof}

For clarity, consider the resulting odd companions.
Put $\overline{\mathcal R}=
(\mathcal R/\mathcal R_{\Lambda\text{-tor}})\otimes L$.
Its quotient by $\mathcal S_L$ is $(D/D_0)\otimes L$.
Choose top Jordan vectors in this finite quotient,
lift them to $y_i\in\overline{\mathcal R}$, and put
$C_i'=(\gamma-\beta_i)^{r_i}y_i\in\mathcal S_L$,
where $r_i$ is the Jordan block length.
For each fixed $\beta$, their residues in
$\mathcal S_L/(\gamma-\beta)\mathcal S_L$ are independent:
a nonzero linear combination divisible by $\gamma-\beta$
would, after subtracting a source preimage, give a nonzero
$\beta$-eigenvector in the torsionfree module
$\overline{\mathcal R}$. Its projection to the independent
socle vectors of the Jordan blocks would be nonzero.
Proposition~\ref{prop:degree-three-saturation} transports
this independence to the three-position coefficient
residues. Thus these odd companions have nonzero
total-product residues in the stated coefficient map.
